% ============================================================================
% 2.7 编译、编译脚本与调试
% ============================================================================
\subsection{clang：汇编成可执行文件}

整个项目只需要一条命令把所有 \texttt{.s} 源文件编译成一个可执行文件：

\begin{lstlisting}[style=bashstyle,caption=ISA/run.sh — 一键构建]
#!/bin/zsh

echo ""
echo "  Building hello.s ..."
echo ""

clang -g -o bin/hello hello.s src/*.s -nostdlib -lSystem

# 然后进入调试器
lldb ./bin/hello
\end{lstlisting}

各参数含义：

\begin{itemize}
    \item \texttt{-g}：在生成的机器码里插入调试信息（DWARF），
        让 lldb 能显示源文件行号
    \item \texttt{-nostdlib}：不链接默认的 C 标准库。我们自己管理内存
    \item \texttt{-lSystem}：但要链接系统调用库（\texttt{libSystem}）。
        macOS 的 \texttt{svc \#0x80} 机制在这里面实现
    \item \texttt{hello.s src/*.s}：所有源文件在一条命令里处理
    \item \texttt{lldb ./bin/hello}：启动后立即进入调试器
\end{itemize}

\subsection{为什么不用 ld 不用自己链接}

一个典型的"手写汇编程序"在 macOS 上其实只需要两件事：
汇编和系统调用。\textbf{\$ clang} 是一个足够智能的汇编器：

\begin{itemize}
    \item 它自动把你写的 \texttt{.s} 文件汇编成机器码
    \item 自动调用链接器把它们拼起来
    \item 自动处理系统调用库引用
\end{itemize}

如果用 \texttt{ld} 单独链接，你需要手动指定 \texttt{-lSystem} 等库路径，
还要保证 \textbf{入口点} 符合约定。统一交给 clang 处理就省掉了。

\subsection{调试流程}

程序在 LLDB 里，常用命令：

\begin{itemize}
    \item \textbf{b store\_string\_ptr\_len} 或 \textbf{b find\_two\_spaces}：在关键函数上下断点
    \item \textbf{r} 或 \textbf{run}：开始运行
    \item \textbf{si}：单步执行一条汇编指令
    \item \textbf{n}：执行下一行（不进入子函数）
    \item \textbf{register read x0}：读寄存器
    \item \textbf{memory read --size 8 --format x --count 16 \$x1}：查看内存内容
\end{itemize}

要验证分配器是否正确，可以在 \texttt{store\_string\_ptr\_len} 上下断点，
看 \texttt{x0}（字符串地址）、\texttt{x1}（序号）、\texttt{x2}（长度）是否
都正确写入 \texttt{string\_ptrs[]} 和 \texttt{string\_lens[]} 中。

\subsection{运行流程}

总结一下：

\begin{verbatim}
./run.sh
  → clang 汇编、链接所有 .s
  → lldb 启动调试器
  → 在程序里输入任意文本
  → 读入 → 复制 → 保存 → 打印 → 分析空格
\end{verbatim}

这个过程中，你可以在任何地方打断点调试器看寄存器的值，看内存的内容，
观察函数调用栈，确认整个链路的每个环节的数据流动方向。
